\section{Conclusions}
We have introduced a novel combination of width-based planning with learning techniques. The learning employs a VAE to learn relevant features in Atari games given images of screen states as training data. The planning is done with RolloutIW(1) using the features learned by the VAE. 
Our approach reduces the size of the feature set from the 20.5 million B-PROST features used in previous work in connection with RolloutIW, to only 4,500. Our algorithm, VAE-IW, outperforms the previous methods, proving that VAEs can learn meaningful representations that can be effectively used for planning with RolloutIW.

In VAE-IW, the symbolic representations are learned from data collected by RolloutIW using B-PROST features. Increasing the diversity and quality of the training data could potentially lead to better representations, and thus better planning results.
One possible way to achieve this could be to iteratively retrain and fine-tune the VAEs based on data collected by the current iteration of VAE-IW. 
The quality of the representations could also be improved by using more expressive discrete models, for example with a hierarchy of discrete latent variables \cite{van2017neural,razavi2019generating}.
Similarly, in the ablation study where we compare this approach with quantizing continuous representations, more flexible models~\cite{ranganath2016hierarchical,sonderby2016ladder,vahdat2020nvae} and more advanced quantization strategies could be employed.
Finally, we can expect further improvements orthogonal to this work, by learning a rollout policy for more effective action selection as in \citet{RLIW}.
